% =====================================================================
%  6. Monitoramento, avaliação e indicadores
% =====================================================================
\section{Monitoramento, avaliação e indicadores}
\label{sec:monitoramento}

\subsection{Indicadores da ação}

O desempenho do Subcomponente 3.2.b será acompanhado pelos indicadores do
\autoref{qua:indicadores}, que apresenta a fórmula de cálculo, o horizonte
temporal, a unidade de medida e a meta final de cada um.

\begin{quadro}[htbp]
\caption{Indicadores de implantação e reabilitação dos sistemas de abastecimento de água e esgotamento sanitário no meio rural}
\label{qua:indicadores}
\footnotesize
\renewcommand{\arraystretch}{1.25}
\begin{tabularx}{\textwidth}{|L|L|Q{1.6cm}|P{2.2cm}|Q{1.4cm}|}
\hline
\cab{Indicador} & \cab{Fórmula de cálculo} & \cab{Horizonte} &
\cab{Unidade} & \cab{Meta final} \\ \hline
Projetos de poços tubulares com sistemas de abastecimento de água implantados &
Quantidade de projetos executados & 1º ao 6º ano & Projetos (nº) & 83 \\ \hline
Domicílios atendidos com sistemas de abastecimento de água &
Soma das comunidades dentro dos municípios contemplados pelas contratações &
1º ao 6º ano & Domicílios & 5.810 \\ \hline
Domicílios atendidos com esgotamento sanitário &
70\% dos domicílios atendidos com abastecimento de água &
4º ao 6º ano & Domicílios & 4.067 \\ \hline
Domicílios atendidos com esgotamento sanitário nas propriedades rurais do IDR-Paraná &
Obtidos a partir dos dados do Censo 2022 & 4º ao 6º ano & Domicílios & 2.300 \\ \hline
População rural atendida com acesso à água potável e ao esgotamento sanitário &
Soma dos habitantes beneficiados pelas obras & 4º ao 6º ano &
Pessoas (nº) & 22.708 \\ \hline
\end{tabularx}
\fonte{Adaptado de \citeonline[Quadro 27]{mop2026}.}
\end{quadro}

\pendencia{Revisar as fórmulas do \autoref{qua:indicadores}: (i) o indicador de
domicílios com abastecimento é calculado como ``soma das comunidades''
(unidade diferente de domicílios); (ii) o MOP escreve ``70\% do total de
pessoas atendidas domicílios atendidos com abastecimento'', o que mistura
pessoas e domicílios; (iii) o horizonte do esgotamento (4º ao 6º ano)
não coincide com o prazo das \atividade{3.2.2.4} e \atividade{3.2.2.5} (anos 2
a 5). O MOP também remete a um ``Quadro XX'' que deve ser corrigido.}

\subsection{Contribuição para os indicadores de resultado do Componente 3}

O Subcomponente 3.2.b contribui para três indicadores de resultado do
Componente~3, cujas metas progressivas estão na
\autoref{tab:metas-componente3}. Os valores de pessoas atendidas consideram
2,8 moradores por domicílio \cite{ibge2022} e incluem, no caso do
abastecimento de água, os sistemas implantados pela Sanepar.

\begin{table}[htbp]
\centering
\caption{Metas progressivas dos indicadores do Componente 3 relacionados ao Subcomponente 3.2.b}
\label{tab:metas-componente3}
\footnotesize
\renewcommand{\arraystretch}{1.2}
\begin{tabular}{lrrrrrrr}
\toprule
\textbf{Indicador} & \textbf{2025} & \textbf{2027} & \textbf{2028} &
\textbf{2029} & \textbf{2030} & \textbf{2031} & \textbf{2032} \\
\midrule
Pessoas com água gerida de forma segura & 0 & 980 & 10.900 & 23.000 & 35.400 & 49.000 & 59.780 \\
\quad sexo feminino & 0 & 400 & 5.200 & 11.000 & 16.000 & 23.300 & 28.500 \\
\quad jovens (15 a 24 anos) & 0 & 110 & 1.300 & 2.800 & 4.200 & 5.900 & 7.000 \\
Pessoas com, pelo menos, esgotamento sanitário & 0 & 126 & 3.000 & 6.700 & 10.600 & 14.600 & 17.800 \\
\quad sexo feminino & 0 & 60 & 1.400 & 3.200 & 5.000 & 6.900 & 8.000 \\
\quad jovens (15 a 24 anos) & 0 & 15 & 360 & 800 & 1.200 & 1.700 & 2.000 \\
Associações RWSS com modelo de gestão adotado & 0 & 0 & 0 & 10 & 30 & 90 & 122 \\
\bottomrule
\end{tabular}
\fonte{Adaptado de \citeonline[Quadro 23]{mop2026}.}
\end{table}

\subsection{Procedimentos de monitoramento e validação}

O monitoramento seguirá os procedimentos definidos nos quadros de etapas das
atividades (\autoref{qua:etapas-3221} e \autoref{qua:etapas-3223}):

\begin{enumerate}[label=\alph*)]
  \item \textbf{supervisão técnica contínua} das obras e da implantação da
        gestão, registrada em relatórios de supervisão, laudos de vistoria e
        testes de estanqueidade, validados pelo IAT;
  \item \textbf{avaliação semestral de desempenho}, consolidada em relatórios
        de conformidade técnica e ambiental, com dados lançados no SIGARH e no
        GeoPR e validação pela UGP à luz do MGAS;
  \item \textbf{reporte anual dos indicadores de resultado} à UGP pelos
        executores (IAT e Sanepar), com base em evidências técnicas de
        implementação;
  \item \textbf{encerramento}, com relatório final de gerenciamento e obras
        concluídas e emissão do Termo de Recebimento Definitivo (TRD),
        validados pela UGP e pelo Banco Mundial.
\end{enumerate}
